\section{Sistema de variables y operacionalización}

\subsection{Definición conceptual de las variables}
La variable independiente es la configuración del detector: arquitectura
(YOLOv8n o YOLOv11) y cuantización (FP32 o INT8). La variable dependiente es el
desempeño del detector: precisión (incluida la robustez) y eficiencia. La
oportunidad de la retroalimentación, la carga cognitiva
\parencite{hart1988} y la usabilidad \parencite{brooke1996} son variables
secundarias.

\subsection{Operacionalización}
La Tabla~\ref{tab:operacionalizacion} resume las dimensiones, los indicadores y
la escala de medición.

\begin{table}[htbp]
  \centering
  \caption{Operacionalización de las variables}
  \label{tab:operacionalizacion}
  \small
  \renewcommand{\arraystretch}{1.15}
  \begin{tabularx}{\textwidth}{@{}>{\raggedright\arraybackslash}p{2.6cm}>{\raggedright\arraybackslash}p{2.6cm}>{\raggedright\arraybackslash}X>{\raggedright\arraybackslash}p{2.0cm}@{}}
    \toprule
    \textbf{Variable} & \textbf{Dimensión} & \textbf{Indicadores} & \textbf{Escala} \\
    \midrule
    VI: Arquitectura de detección & Arquitectura & YOLOv8n; YOLOv11 & Nominal \\
     & Cuantización & FP32; INT8 & Nominal \\
    \midrule
    VD: Desempeño del detector & Precisión (incl. robustez) & mAP@0.5; mAP@0.5:0.95; F1 macro; Kappa; AUC-ROC macro; CER; WER & Razón / intervalo \\
     & Eficiencia & Latencia (ms); memoria (MB); tamaño (MB); $\Delta$mAP & Razón \\
    \midrule
    Secundaria: Retroalimentación docente & Oportunidad & Tiempo de revisión por página (min) & Razón \\
    Secundaria: Carga cognitiva & Carga global & NASA-TLX (0--100) & Intervalo \\
    Secundaria: Usabilidad & Usabilidad global & SUS (0--100) & Intervalo \\
    \bottomrule
  \end{tabularx}
  \notaTabla{Elaboración propia a partir de la literatura revisada.}
\end{table}
